\input{sol_pre.tex}
\input{q_2025-10-05_1.tex}
\begin{document}
\input{sol_head.tex}

\Savol{1}
$x$, $y$ ikkalasi toq bölsa $x+y$ juft; aks holda $x$ yoki $y$ juft. Demak $(x+y)xy$ har doim juft. ($x=1$, $y=2$ da A, B, C toq.)
\Javob{D}

\Savol{2}
$\{-1{,}6\}=0{,}4$, \ $[-1{,}6]=-2$, \ $[6{,}6]=6$, \ $\{6{,}6\}=0{,}6$:
\[\left\{\frac{0{,}4}{-2}-\frac6{0{,}6}\right\}=\{-0{,}2-10\}=\{-10{,}2\}=0{,}8=\frac45.\]
\Javob{B}

\Savol{3}
$n$ qator, har birida $n+4$ örindiq: $n(n+4)=60$ $\Rightarrow$ $n=6$.
\Javob{C}

\Savol{4}
Çegirma: $500\,000-410\,000=90\,000$; \ $\dfrac{90\,000}{500\,000}=0{,}18=18\%$.
\Javob{A}

\Savol{5}
Surat: $x^{\frac32-\frac15}=x^{\frac{13}{10}}$; \ maxraj: $\left(x^4\cdot x^{\frac52}\right)^{\frac15}=x^{\frac{13}{10}}$. Natija $1$.
\Javob{A}

\Savol{6}
$\sqrt{3\pm2\sqrt2}=\sqrt2\pm1$:
\[\frac{(\sqrt2+1)-(\sqrt2-1)}{(\sqrt2+1)+(\sqrt2-1)}=\frac2{2\sqrt2}=\frac{\sqrt2}2.\]
\Javob{C}

\Savol{7}
$\sqrt{10\pm2\sqrt{21}}=\sqrt7\pm\sqrt3$: \ $\left(2\sqrt7\right)^2=28$.
\Javob{D}

\Savol{8}
$S_9=9a_5=54$ $\Rightarrow$ $a_5=6$; \ $a_6^2=100-36=64$, kamayuvçi: $a_6=-8$.\\
$d=-14$, \ $a_4=a_5-d=20$.
\Javob{C}

\Savol{9}
$b=aq$, $c=aq^2$: \ $2\cdot1{,}5b=a+2c$ $\Rightarrow$ $3q=1+2q^2$ $\Rightarrow$ $q=1$ yoki $q=\frac12$. Kamayuvçi: $q=\frac12$.
\Javob{D}

\Savol{10}
$x^3=2$: \ $y^3x^{-3}=\frac82=4$; \ $x^{-1}\cdot y\cdot x^{-5}\cdot x^3=yx^{-3}=1$; \ maxraj $x^{-3}(y+1)=\frac32$.
\[\frac{4-1}{3/2}=2.\]
\Javob{C}

\Savol{11}
$u=\sqrt a$, $v=\sqrt b$:
\[\frac{u-v}{u(u+v)}+\frac{u+v}{u(u-v)}-\frac{2v^2}{u(u-v)(u+v)}=\frac{(u-v)^2+(u+v)^2-2v^2}{u\left(u^2-v^2\right)}=\frac{2u}{u^2-v^2}.\]
\[\frac{2u}{(u-v)(u+v)}\cdot\frac{v-u}{uv}=\frac{-2}{v(u+v)}=\frac{-2}{\sqrt{ab}+b}.\]
\Javob{C}

\Savol{12}
Çap tomon $\pm1$: \ $x=1$ ($\sin1>0$ — mos), \ $x=-1$ ($\sin(-1)<0$ — mos). $2$ ta yeçim.
\Javob{B}

\Savol{13}
$\tg t+\ctg t=\dfrac2{\sin2t}$:
\[\frac2{\sin4\alpha}+\frac2{\sin12\alpha}=\frac{2(\sin12\alpha+\sin4\alpha)}{\sin4\alpha\sin12\alpha}=\frac{4\sin8\alpha\cos4\alpha}{\sin4\alpha\sin12\alpha}=\frac{8\cos^24\alpha}{\sin12\alpha}.\]
\Javob{D}

\Savol{14}
$t=2^{\sqrt x}\ge1$: \ $t^2-5t-24<0$ $\Rightarrow$ $t<8$ $\Rightarrow$ $\sqrt x<3$ $\Rightarrow$ $0\le x<9$. Butunlar: $0,1,\dots,8$ — $9$ ta.
\Javob{B}

\Savol{15}
$x>0$, $t=\log_6x$: \ $5t^2<3+2t$ $\Rightarrow$ $(5t+3)(t-1)<0$ $\Rightarrow$ $-\frac35<t<1$ $\Rightarrow$ $6^{-0{,}6}<x<6$.\\
Butunlar $1,\dots,5$: \ $1+5=6$.
\Javob{B}

\Savol{16}
$t=x^3$: $t^2+\sqrt{2021}\,t-1=0$ — ildizlari köpaytmasi $-1$, ya'ni ikkita haqiqiy ildiz; har biriga bitta $x$. Jami $2$ ta.
\Javob{C}

\Savol{17}
$(2x-4)^2(x+2)^2=4\left(x^2-4\right)^2$. \ $u=x^2-4$: \ $4u^2-u-3=0$ $\Rightarrow$ $u=1$ yoki $u=-\frac34$.\\
$x^2=5$ yoki $x^2=\frac{13}4$ — jami $4$ ta ildiz.
\Javob{D}

\Savol{18}
\[\frac{x^3+1-2x^2+3x-1}{(2x-1)(x-1)}>0\ \Rightarrow\ \frac{x\left(x^2-2x+3\right)}{(2x-1)(x-1)}>0.\]
$x^2-2x+3>0$, demak $\dfrac x{(2x-1)(x-1)}>0$: \ $x\in\left(0;\frac12\right)\cup(1;\infty)$.
\Javob{A}

\Savol{19}
Aniqlaniş sohasi $x\ge-1$. $\sqrt{x+1}=0$ da ($x=-1$) bajariladi, $x>-1$ da $x\le0$. Yeçim $[-1;0]$: butunlar $-1$, $0$.
\Javob{A}

\Savol{20}
Tarmoqlar pastga: $a<0$. Uçi $Oy$ dan öngda: $-\frac b{2a}>0$ $\Rightarrow$ $b>0$. $Oy$ ni musbat qismida kesadi: $c>0$.
\Javob{C}

\Savol{21}
\[\int_0^2xe^{x^2}dx=\frac12e^{x^2}\Big|_0^2=\frac12\left(e^4-1\right).\]
\Javob{B}

\Savol{22}
$y'=2x+4=3$ $\Rightarrow$ $x_0=-\frac12$, \ $y_0=\frac14-2-2=-3{,}75$.
\Javob{A}

\Savol{23}
$f=\left(x^2+3x\right)\left(x^2+3x+2\right)$, \ $t=x^2+3x\ge-\frac94$: \ $t(t+2)=(t+1)^2-1\ge-1$ ($t=-1$ da).
\Javob{C}

\Savol{24}
Urinma va kesuvçi: $AB^2=AD\cdot AE$ $\Rightarrow$ $576=18\cdot AE$, $AE=32$.\\
$DE=14$ — diametr, $r=7$: \ $AO=18+7=25$.
\Javob{D}

\Savol{25}
$ABE$ uçburçak va $F,O,C$ töğri çiziq uçun Menelay teoremasi:
\[\frac{AF}{FB}\cdot\frac{BC}{CE}\cdot\frac{EO}{OA}=1\ \Rightarrow\ 3\cdot\frac43\cdot\frac{EO}{OA}=1\ \Rightarrow\ \frac{EO}{OA}=\frac14.\]
$BCF$ uçburçak va $A,O,E$ uçun: \ $\frac{BA}{AF}\cdot\frac{FO}{OC}\cdot\frac{CE}{EB}=1$ $\Rightarrow$ $\frac43\cdot\frac{FO}{OC}\cdot3=1$ $\Rightarrow$ $\frac{FO}{OC}=\frac14$.\\
$FOE$ va $COA$ uçburçaklarda $O$ dagi burçaklar teng (vertikal):
\[\frac{S_{FOE}}{S_{AOC}}=\frac{FO\cdot OE}{OC\cdot OA}=\frac1{16}\ \Rightarrow\ S_{FOE}=\frac{144}{16}=9.\]
\Javob{B}

\Savol{26}
Bissektrisa: $\frac{AD}{DB}=\frac{AC}{BC}=\frac34$ $\Rightarrow$ $AD=\frac{15}7$. \ $AM=\frac52$:
\[MD=\frac52-\frac{15}7=\frac5{14}.\]
\Javob{B}

\Savol{27}
$BO:OD=1:2$ $\Rightarrow$ $BC:AD=1:2$. $S_{BOC}=s$ bölsa, $S_{AOB}=S_{COD}=2s$, $S_{AOD}=4s$:
\[9s=162\ \Rightarrow\ s=18,\qquad S_{ABO}=36.\]
\Javob{B}

\Savol{28}
Muntazam beşburçakda har bir burçak $108^\circ$ ni diagonallar $36^\circ$ li uçta teng qismga böladi. $AC$ diagonalda $AF=GC=x$, $FG=2$, $AB=a$.\\
$ABF$: $\angle BAF=\angle ABF=36^\circ$ $\Rightarrow$ $BF=AF=x$.\\
$ABG$: $\angle BAG=36^\circ$, $\angle ABG=72^\circ$ $\Rightarrow$ $\angle AGB=72^\circ$, $AB=AG=x+2$.\\
$BFG$ va $ABG$ — ikkalasining burçaklari $36^\circ,72^\circ,72^\circ$, öxşaş:
\[\frac{FG}{BG}=\frac{BG}{AB}\ \Rightarrow\ x^2=2(x+2)\ \Rightarrow\ x=1+\sqrt5,\quad a=3+\sqrt5.\]
Perimetr: $5\left(3+\sqrt5\right)$.
\Javob{C}

\Savol{29}
$ABCD$ parallelogramm: $\vec{AD}=\vec{BC}$ $\Rightarrow$ $D(4;\,-2)$.\\
$\vec{AC}=(12;\,9)$, $|AC|=15$; \ $\vec{DB}=(-6;\,7)$, $|DB|=\sqrt{85}$:
\[\cos\varphi=\frac{-72+63}{15\sqrt{85}}=-\frac3{5\sqrt{85}}.\]
\Javob{A}

\Savol{30}
$CH\perp AB$: \ $AB=25$, $CH=\frac{15\cdot20}{25}=12$. Uç perpendikulyar teoremasiga köra $DH\perp AB$:
\[DH=\sqrt{35^2+12^2}=37.\]
\Javob{B}

\Savol{31}
Mayonez yoki piyozli: $8+5-2=11$. Ikkalasi ham yöq: $14-11=3$.
\Javob{B}

\Savol{32}
$7$ birinçi xonada: $9\cdot9=81$; ikkinçi yoki uçinçi xonada: $8\cdot9=72$ dan. Jami $81+72+72=225$.
\Javob{B}

\newpage
\Savol{33}
Şar radiusi $R$.\\
\textbf{Konus:} balandligi $h$, $r^2=h(2R-h)$: \ $V=\frac\pi3h^2(2R-h)$, \ $V'=\frac\pi3(4Rh-3h^2)=0$ $\Rightarrow$ $h=\frac{4R}3$, \ $R_k^2=\frac{8R^2}9$, \ $R_k=\frac{2\sqrt2}3R$.\\
\textbf{Silindr:} balandligi $H$, $r^2=R^2-\frac{H^2}4$: \ $V=\pi\left(R^2-\frac{H^2}4\right)H$, \ $V'=0$ $\Rightarrow$ $H=\frac{2R}{\sqrt3}$, \ $R_s^2=\frac{2R^2}3$, \ $R_s=\frac{\sqrt2}{\sqrt3}R$.

\textbf{33.} $\dfrac{R_k}{R_s}=\dfrac{2\sqrt2}3\cdot\dfrac{\sqrt3}{\sqrt2}=\dfrac{2\sqrt3}3$ — \textbf{D}.

\textbf{34.} $V_k=\frac13\pi\cdot\frac{8R^2}9\cdot\frac{4R}3=\frac{32\pi R^3}{81}$, \ $V_s=\pi\cdot\frac{2R^2}3\cdot\frac{2R}{\sqrt3}=\frac{4\pi R^3}{3\sqrt3}$:
\[\frac{V_k}{V_s}=\frac{32\cdot3\sqrt3}{81\cdot4}=\frac{8\sqrt3}{27}\ \text{— \textbf{A}}.\]

\textbf{35.} Konus yasovçisi $l=\sqrt{\frac{8}9+\frac{16}9}\,R=\frac{2\sqrt6}3R$: \ $S_k=\pi R_kl=\frac{8\sqrt3}9\pi R^2$; \ $S_s=2\pi R_sH=\frac{4\sqrt2}3\pi R^2$:
\[\frac{S_k}{S_s}=\frac{8\sqrt3}9\cdot\frac3{4\sqrt2}=\frac{2\sqrt3}{3\sqrt2}=\frac{\sqrt6}3\ \text{— \textbf{E}}.\]
\Javob{33 — D, \ 34 — A, \ 35 — E}

\newpage
\Savol{36}
$s=\sqrt x\ge0$, \ $y=\sqrt{9-s}\ge0$. U holda $y^2=9-s$, tenglama esa $s^2=9-y$:
\[\begin{cases}s^2+y=9,\\ y^2+s=9.\end{cases}\]
Ayiramiz: $s^2-y^2-(s-y)=0$ $\Rightarrow$ $(s-y)(s+y-1)=0$.

1) $s=y$: \ $s^2+s-9=0$ $\Rightarrow$ $s=\dfrac{\sqrt{37}-1}2>0$; \ $x=s^2=9-s=\dfrac{19-\sqrt{37}}2$.

2) $s+y=1$: \ $s^2-s-8=0$ $\Rightarrow$ $s=\dfrac{1\pm\sqrt{33}}2$; musbati $s\approx3{,}37$ da $y=1-s<0$ — mos emas.

a) Bitta ildiz. \quad b) $x=\dfrac{19-\sqrt{37}}2$.
\Javob{a) 1; \quad b) $\frac{19-\sqrt{37}}2$}

\newpage
\Savol{37}
$\sin x+\sin3x=2\sin2x\cos x$ $\Rightarrow$ çap tomon $2\sin2x\,(1+\cos x)$.\\
$1+\cos2x+2\cos x=2\cos^2x+2\cos x=2\cos x\,(1+\cos x)$; \ $1+\cos x\ge0$, şuning uçun öng tomon $2|\cos x|(1+\cos x)$:
\[(1+\cos x)\left(\sin2x-|\cos x|\right)=0.\]
1) $\cos x=-1$: \ $x=\pi+2\pi k$.\\
2) $2\sin x\cos x=|\cos x|$: \ $\cos x=0$ $\Rightarrow$ $x=\frac\pi2+\pi k$; \ $\cos x>0$ da $\sin x=\frac12$ $\Rightarrow$ $x=\frac\pi6+2\pi k$; \ $\cos x<0$ da $\sin x=-\frac12$ $\Rightarrow$ $x=\frac{7\pi}6+2\pi k$.

a) $[-2\pi;2\pi]$: \ $\pm\pi$; \ $\pm\frac\pi2,\ \pm\frac{3\pi}2$; \ $\frac\pi6,\ -\frac{11\pi}6$; \ $\frac{7\pi}6,\ -\frac{5\pi}6$ — jami $10$ ta.

b) $[0;2\pi]$: \ $\frac\pi6+\frac\pi2+\pi+\frac{7\pi}6+\frac{3\pi}2=\frac{13\pi}3$.
\Javob{a) 10; \quad b) $\frac{13\pi}3$}

\newpage
\Savol{38}
a) $x=1$: \ $x-\frac2x=-1$: \ $4f(1)+2f(-1)=2$.\\
\phantom{a) }$x=-1$: \ $x-\frac2x=1$: \ $4f(-1)+2f(1)=0$.\\
Ayiramiz: $2\left(f(1)-f(-1)\right)=2$ $\Rightarrow$ $f(1)-f(-1)=1$.

b) $x=2$: \ $x-\frac2x=1$: \ $4f(2)+4f(1)=9$.\\
\phantom{b) }$x=-2$: \ $x-\frac2x=-1$: \ $4f(-2)+4f(-1)=-7$.\\
Ayiramiz: $4\left(f(2)-f(-2)\right)+4\cdot1=16$ $\Rightarrow$ $f(2)-f(-2)=3$.
\Javob{a) 1; \quad b) 3}

\newpage
\Savol{39}
Umumiy nuqtalar: $f(x)=g(x)$ $\Leftrightarrow$
\[a=-x^3+2x^2+5x-6{,}5-0{,}5(x-1)^2=-x^3+1{,}5x^2+6x-7=\varphi(x).\]
Uriniş nuqtasida $f'=g'$, ya'ni $\varphi'(x)=0$ (karrali ildiz):
\[\varphi'(x)=-3x^2+3x+6=-3(x-2)(x+1)=0\ \Rightarrow\ x=2\ \text{yoki}\ x=-1.\]
$\varphi(2)=-8+6+12-7=3$, \ $\varphi(-1)=1+1{,}5-6-7=-10{,}5$. Ikkala holatda ham umumiy nuqtalar ikkita (bittasida urinadi).

a) $a_{\max}=3$.

b) Uriniş nuqtasi $x=2$, \ $y=g(2)=0{,}5+3=3{,}5$: \ $2+3{,}5=5{,}5$.
\Javob{a) 3; \quad b) $5{,}5$}

\newpage
\Savol{40}
Kesişiş nuqtalari $g$ ning öng tarmoğida yotadi: $y=x+k$. $(6;10)$ dan ötadi: $k=4$, demak $Oy$ dagi umumiy nuqta $(0;4)$, $c=4$.\\
Simmetriya öqi $x=4$: \ $b=-8a$, \ $f(x)=ax^2-8ax+4$. \ $f(6)=36a-48a+4=10$ $\Rightarrow$ $a=-\frac12$:
\[f(x)=-\frac{x^2}2+4x+4.\]
a) $f_{\max}=f(4)=-8+16+4=12$.

b) Soha $x\in[0;6]$ da parabola va $y=x+4$ orasida:
\[S=\int_0^6\left(-\frac{x^2}2+4x+4-x-4\right)dx=\int_0^6\left(-\frac{x^2}2+3x\right)dx=-36+54=18.\]
\Javob{a) 12; \quad b) 18}

\newpage
\Savol{41}
\fig{\begin{tikzpicture}[scale=0.55,font=\small]
\coordinate (B) at (-4,0); \coordinate (H) at (0,0); \coordinate (C) at (9,0); \coordinate (A) at (0,6); \coordinate (M) at (-3.328,1.008); \coordinate (N) at (4.992,2.672);
\coordinate (O1) at (-1.394,1.394); \coordinate (O2) at (2.092,2.092); \coordinate (P) at (0,1.672);
\fill[okgreen!12] (A)--(M)--(N)--cycle;
\draw (A)--(B)--(C)--cycle; \draw (A)--(H); \draw (O1) circle (1.394); \draw (O2) circle (2.092); \draw (M)--(N);
\draw[ax] (H)--(O1); \draw[ax] (H)--(O2); \draw[ax] (A)--(O1); \draw[ax] (A)--(O2);
\fill (O1) circle (2.5pt); \fill (O2) circle (2.5pt);
\draw (0,0.35)--(0.35,0.35)--(0.35,0);
\node[above] at (A) {$A$}; \node[below left] at (B) {$B$}; \node[below right] at (C) {$C$}; \node[below] at (H) {$H$}; \node[left] at (M) {$M$}; \node[above right] at (N) {$N$};
\node[below left] at (O1) {$O_1$}; \node[below right] at (O2) {$O_2$}; \node[aux,above left] at (P) {$P$};
\node[below] at (-2,0) {$4$}; \node[below] at (4.5,0) {$9$};
\end{tikzpicture}}
$AH^2=BH\cdot HC=36$, \ $AH=6$. $\angle A=90^\circ$; $\triangle ABH\sim\triangle CAH$, şuning uçun içki aylanalar radiuslari nisbati $\frac{r_1}{r_2}=\frac{AB}{AC}$.\\
$HO_1$, $HO_2$ — töğri burçaklar bissektrisalari: $\angle O_1HO_2=90^\circ$, \ $HO_1=r_1\sqrt2$, $HO_2=r_2\sqrt2$, ya'ni $\frac{HO_1}{HO_2}=\frac{AB}{AC}$.\\
Demak $\triangle O_1HO_2\sim\triangle BAC$ (katetlar böyiça) va $\angle HO_1O_2=\angle B$.\\
$MN$ va $AH$ kesişiş nuqtasi $P$: \ $\triangle HO_1P$ da $\angle PHO_1=45^\circ$, $\angle HO_1P=\angle B$ $\Rightarrow$ $\angle APM=45^\circ+\angle B$.\\
$\triangle APM$: \ $\angle PAM=90^\circ-\angle B$ $\Rightarrow$ $\angle AMN=180^\circ-(90^\circ-\angle B)-(45^\circ+\angle B)=45^\circ$.

a) $\angle MAN=90^\circ$, $\angle AMN=45^\circ$ $\Rightarrow$ $\alpha=\angle ANM=45^\circ$, \ $\tg\alpha=1$.

b) $\triangle AMO_1=\triangle AHO_1$: $AO_1$ umumiy, $\angle MAO_1=\angle HAO_1$ ($AO_1$ — bissektrisa), $\angle AMO_1=\angle AHO_1=45^\circ$. Demak $AM=AH=6$, xuddi şunday $AN=AH=6$:
\[S_{AMN}=\frac12\cdot6\cdot6=18.\]
\Javob{a) 1; \quad b) 18}

\newpage
\Savol{42}
\fig{\begin{tikzpicture}[scale=0.5,font=\small]
\coordinate (A) at (0,0); \coordinate (B) at (8,0); \coordinate (C) at (8,8); \coordinate (D) at (0,8);
\fill[okgreen!12] (0,0) rectangle (4,4); \fill[okgreen!12] (5,5) rectangle (8,8);
\draw (A)--(B)--(C)--(D)--cycle; \draw (4,0)--(4,4)--(0,4); \draw (5,8)--(5,5)--(8,5);
\draw[ax] (A)--(C);
\node[below left] at (A) {$A$}; \node[below right] at (B) {$B$}; \node[above right] at (C) {$C$}; \node[above left] at (D) {$D$};
\node[below] at (4,0) {$G$}; \node[left] at (0,4) {$E$}; \node[below right] at (4,4) {$F$}; \node[above left] at (5,5) {$K$}; \node[right] at (8,5) {$M$}; \node[above] at (5,8) {$L$};
\node[below] at (2,0) {$4$}; \node[above] at (6.5,8) {$3$};
\end{tikzpicture}}
Kvadratlar tomonlari $4$ va $3$. $F$ va $K$ — $A$ va $C$ burçaklardagi kvadratlarning qarama-qarşi uçlari, ikkalasi $AC$ diagonalda yotadi:
\[AF=4\sqrt2,\quad CK=3\sqrt2,\quad AC=a\sqrt2\ \ (a=AB).\]
$FK=|AC-AF-CK|=\sqrt2\,|a-7|=\sqrt2$ $\Rightarrow$ $a=8$ yoki $a=6$.

a) $S_{\max}=8^2=64$.

b) $a=6$: \ $ED=6-4=2$, $FE=4$, $\angle FED=90^\circ$:
\[FD=\sqrt{4^2+2^2}=2\sqrt5.\]
\Javob{a) 64; \quad b) $2\sqrt5$}

\newpage
\Savol{43}
\fig{\begin{tikzpicture}[scale=0.8,font=\small]
\coordinate (A) at (0,0); \coordinate (B) at (4,0); \coordinate (C) at (4,4); \coordinate (D) at (0,4); \coordinate (E) at (2,3.464); \coordinate (O) at (2.536,2.536);
\fill[okgreen!12] (O)--(E)--(C)--cycle;
\draw (A)--(B)--(C)--(D)--cycle; \draw (A)--(C); \draw (B)--(E); \draw (A)--(E)--(C);
\draw[ax] (E)--(D);
\node[below left] at (A) {$A$}; \node[below right] at (B) {$B$}; \node[above right] at (C) {$C$}; \node[above left] at (D) {$D$};
\node[above] at (E) {$E$}; \node[right] at (O) {$O$};
\end{tikzpicture}}
a) Kvadrat içida $ABE'$ teng tomonli uçburçak yasaymiz. $AE'=AB=AD$, $\angle E'AD=30^\circ$ $\Rightarrow$ $\angle ADE'=75^\circ$. Xuddi şunday $BE'=BC$, $\angle E'BC=30^\circ$ $\Rightarrow$ $\angle BCE'=75^\circ$, \ $\angle E'CD=15^\circ$.\\
Şuningdek $\angle E'AC=60^\circ-45^\circ=15^\circ$. Demak $E'$ ikkala şartni qanoatlantiradi; bunday nuqta bitta ($A$ va $C$ dan çiqqan nurlar kesişmasi), ya'ni $E=E'$:
\[\angle ABE=60^\circ.\]
b) $\angle ACE=\angle BCE-\angle BCA=75^\circ-45^\circ=30^\circ$. $\triangle ABO$ da $\angle BAO=45^\circ$, $\angle ABO=60^\circ$ $\Rightarrow$ $\angle AOB=75^\circ$, demak $\angle EOC=75^\circ$.\\
$\triangle OEC$: $\angle OCE=30^\circ$, $\angle EOC=75^\circ$ $\Rightarrow$ $\angle OEC=75^\circ$, ya'ni $CO=CE$. Sinuslar teoremasi:
\[CE=\frac{OE\sin75^\circ}{\sin30^\circ}=2\left(\sqrt3-1\right)\cdot\frac{\sqrt2\left(\sqrt3+1\right)}4=\sqrt2.\]
\[S_{OEC}=\frac12\cdot CO\cdot CE\cdot\sin30^\circ=\frac14\cdot2=\frac12.\]
\Javob{a) $60^\circ$; \quad b) $\frac12$}

\newpage
\Savol{44}
\fig{\begin{tikzpicture}[scale=1.1,font=\small]
\coordinate (S) at (0,2.83); \coordinate (P) at (-1.415,0); \coordinate (Q) at (1.415,0); \coordinate (O) at (0,0);
\coordinate (G) at (-0.586,1.657); \coordinate (K) at (0.586,1.657); \coordinate (G0) at (-0.586,0); \coordinate (K0) at (0.586,0); \coordinate (T) at (0,1.657);
\draw (S)--(P)--(Q)--cycle; \draw[okgreen!70!black,thick] (G0)--(G)--(K)--(K0);
\draw[ax] (S)--(O); \draw[ax] (G)--(K);
\node[above] at (S) {$S$}; \node[left] at (G) {$G$}; \node[right] at (K) {$K$}; \node[below] at (O) {$O$};
\node[below left] at (P) {}; \node[aux,above right] at (T) {$T$};
\end{tikzpicture}}
Asos tomoni $a$, balandligi $2a$; konus radiusi $R$, balandligi $H=2R$. Yuqori asos uçlari ($G$, $K$) konus yasovçilarida; $GK$ — yuqori asos diagonali, öq kesimda yotadi. $TK=\frac{a}{\sqrt2}$, $ST=H-2a$:
\[\frac{TK}{R}=\frac{ST}{H}\ \Rightarrow\ \frac a{\sqrt2}=\frac{2R-2a}2=R-a\ \Rightarrow\ R=a\left(1+\frac1{\sqrt2}\right).\]
Yon sirt: $4\cdot a\cdot2a=8a^2=24-16\sqrt2$ $\Rightarrow$ $a^2=3-2\sqrt2$, \ $a=\sqrt2-1$.\\
$R=\left(\sqrt2-1\right)\cdot\frac{2+\sqrt2}2=\frac{\sqrt2}2$, \ $H=\sqrt2$.

a) $l=\sqrt{\frac12+2}=\frac{\sqrt{10}}2$: \ $S_{\text{yon}}=\pi Rl=\pi\cdot\frac{\sqrt2}2\cdot\frac{\sqrt{10}}2=\frac{\sqrt5\,\pi}2$.

b) $\angle GSK$ — öq kesimning uçidagi burçak, $\tg\frac{\angle GSK}2=\frac RH=\frac12$:
\[\sin\angle GSK=\frac{2\cdot\frac12}{1+\frac14}=\frac45.\]
\Javob{a) $\frac{\sqrt5\,\pi}2$; \quad b) $\frac45$}

\newpage
\Savol{45}
$D$ dan $AB$ ga perpendikulyar $DH=2$, $OD=3$ $\Rightarrow$ $OH=\sqrt{9-4}=\sqrt5$.

a) Parabola öqi $AB$ ga perpendikulyar, şuning uçun $C$ va $D$ simmetrik: $CD=2\cdot OH=2\sqrt5$.

b) $AB$ ni abssissa öqi, $O$ ni boşlanğiç nuqta deb olsak, parabola $y=kx^2$, $D\left(\sqrt5;2\right)$: \ $k=\frac25$. Böyalgan soha — parabola, $AB$ va $DH$ orasida:
\[S=\int_0^{\sqrt5}\frac25x^2dx=\frac25\cdot\frac{5\sqrt5}3=\frac{2\sqrt5}3.\]
\Javob{a) $2\sqrt5$; \quad b) $\frac{2\sqrt5}3$}

\ShowJavoblar
\end{document}
